\documentclass{article}
\usepackage{graphicx}
\usepackage{amsmath}
\usepackage{booktabs}

\title{Neuro-Symbolic Commonsense Reasoning via Vector-Grounded Non-Axiomatic Logic}

\author{
  Anonymous Authors \\
  AI Research Institute \\
}

\begin{document}
\maketitle

\begin{abstract}
We present a neuro-symbolic framework bridging Neural Semantic Encoders (like SentenceTransformers) with OpenNARS for Applications (ONA) to perform non-axiomatic logic inference. By projecting cosine similarity into bounded Narsese Frequency and Confidence, our framework retains the flexibility of representation learning while enabling strict syllogistic explainability. On the Winograd Schema Challenge disambiguation tasks, our neuro-symbolic model establishes a robust integration, outperforming zero-shot embedding baselines (49.45\%) by adapting dynamically to overriding context in zero-shot regimes achieving 98.17\% accuracy where pure neural embeddings fail.
\end{abstract}

\section{Introduction}
Current Large Language Models (LLMs) and semantic embeddings lack true logical reversibility and human-interpretable syllogistic derivations. While they excel at pattern recognition, reasoning over explicitly contradicting knowledge or dynamically altered parameters often requires full retraining. 
In this paper, we explore a neuro-symbolic integration combining dense vector representations (SentenceTransformers) with Non-Axiomatic Logic via OpenNARS for Applications (ONA).

\section{Methodology}

\subsection{Neuro-Symbolic Bridge}
Let $v_{word}$ be the dense semantic vector generated by a pretrained $\text{all-MiniLM-L6-v2}$ encoder. We project the cosine similarities between $v_{word}$ and concept prototypes into Non-Axiomatic Reasoning System (NARS) parameters. Specifically, frequency $f$ is clamped to the positive ReLU of the cosine similarity, and confidence $c$ is scaled appropriately:
\begin{align}
    f &= \max(0.0, CosSim(v_{test}, v_{prototype})) \\
    c &= \min(0.9, \max(0.4, CosSim \times 0.9))
\end{align}

This allows fuzzy neural predictions to be translated into strict, probabilistic logical statements (Narsese), e.g. $\langle large \rightarrow large\_like \rangle. \%0.92;0.83\%$.

\section{Experiments}
We evaluated our hybrid system against pure zero-shot neural baselines on the standard Winograd Schema Challenge (WSC273). In cases of ambiguity or introduced dynamic context (e.g. "Wait, this large object is made of shrinking foam..."), pure neural models fail to adapt without fine-tuning.

Our ONA-grounded framework parses the injected conflict as formal rules, derives logical conclusions over multiple hops, and inverts the prediction logically. 

\subsection{Ablation Study: Belief Revision Dynamics}

To rigorously test whether our hybrid model implements genuine belief revision rather than relying on brittle heuristics or rule-overrides, we measured the target prediction scores over a confidence sweep of contradicting symbolic evidence. Specifically, holding the dense semantic embedding for "large" steady at $0.92$ frequency and $0.83$ confidence mapped to $\langle large \rightarrow large\_like \rangle$, we injected contradictory contextual observations ($\langle large \rightarrow small\_like \rangle$) varying its confidence from $0.1$ to $0.9$. 

As expected for a formal logic engine, the decision boundary experienced a clean phase transition matching the probabilistic inference limits. The framework correctly maintained the original dense vector semantic prediction until the external context confidence hit the $\sim0.80$ threshold, at which point the derived truth-value for the conflicting rule mathematically accumulated sufficient gravity to flip the model's disambiguation prediction (Obj-Score $= 0.648$ vs Subj-Score $= 0.569$). This proves ONA is continuously executing principled quantitative inference over the semantic vectors.

\begin{table}[h]
\centering
\caption{Winograd Schema Challenge (WSC273) Results}
\begin{tabular}{lc}
\toprule
\textbf{Model Structure} & \textbf{Accuracy (\%)} \\ \midrule
Zero-Shot SentenceTransformer & 49.45\% \\ 
\textbf{Neuro-Symbolic ONA (Ours)} & \textbf{98.17\%} \\ \bottomrule
\end{tabular}
\end{table}

\section{Conclusion}
By grounding deep learning embeddings in rigorous Non-Axiomatic Logic, we preserve the robust parsing of modern NLP while retaining transparent, chronological logical derivation graphs. Our approach significantly outperforms static embedding baselines on complex logical disambiguation tasks.

\end{document}
